\documentclass[xcolor=dvipsnames,t,aspectratio=169]{beamer}

\usecolortheme{rose}
\usecolortheme{dolphin}
\usetheme{Boadilla}

\input{../../templates/slides/imports}
\input{../../templates/slides/settings}
\input{../../templates/slides/commands}

\usepackage{tikz}
\usepackage{svg}
\usetikzlibrary{arrows.meta, shapes}

\titlegraphic{
    \includegraphics[scale = 0.5]{../../templates/slides/logo}
}

\logo{
\begin{tikzpicture}[overlay,remember picture]
\node[left=1.1cm, below=0.2cm] at (current page.30){
    \includegraphics[width=0.1\textwidth]{../../templates/slides/logo}};
\end{tikzpicture}
}

\newcommand{\highlight}[1]{{\color{nes_dark_orange} #1}}

\title{Introdução ao Linux, Shell e Organização de Projetos}

\author{Eduardo Adame}

\date{{\color{nes_dark_purple}\textbf{Ciência de Dados}\\[0.5em] 04 de março de 2026}}

\begin{document}

\frame[plain]{\titlepage}
\setcounter{framenumber}{0}

\section{Sobre mim}
\begin{frame}[c]{Sobre mim}
    \begin{columns}[c]
        \begin{column}{0.7\textwidth}
            \begin{itemize}
                \item \textbf{Nome:} Eduardo Adame
                \item \textbf{Formação:}
                \begin{itemize}
                    \item Mestrando em Matemática Aplicada e Ciência de Dados (FGV EMAp); 
                    \item Bacharel em Ciência de Dados e Inteligência Artificial (FGV EMAp).
                \end{itemize}
                \item \textbf{Experiência:} Diversos cursos produzidos/ministrados, ex-consultor por 2 anos e meio no Banco Mundial.
                \item \textbf{Pesquisa:} Estatística bayesiana; quantificação de incerteza, modelos não-paramétricos, inferência exata.
                \item \textbf{Contato:} eadamesalles@gmail.com
            \end{itemize}
        \end{column}
        \begin{column}{0.3\textwidth}
            \centering
            \includegraphics[width=0.9\textwidth]{images/me.jpg}
        \end{column}
    \end{columns}
\end{frame}

\begin{frame}[c]{Sobre o curso}
\begin{display}[Objetivos]
\begin{itemize}
\item Aprender a trabalhar com \highlight{dados reais}
\item Construir \highlight{pipelines de ciência de dados}
\item Dominar ferramentas fundamentais do ecossistema:
\begin{itemize}
\item Linux
\item Git
\item Python
\item Bancos de dados
\end{itemize}
\end{itemize}
\end{display}


\end{frame}


\begin{frame}[c]{Sobre o curso}

\begin{display}[Metodologia]
\begin{itemize}
\item Aulas teóricas curtas
\item Demonstrações ao vivo
\item Exercícios semanais
\item Projetos práticos
\end{itemize}
\end{display}

\end{frame}

\begin{frame}[c]{Estrutura do curso - 15 semanas}
    \begin{columns}[c]
        \begin{column}{0.5\textwidth}
            \textbf{Fundamentos (Semanas 1-8):}
            \begin{enumerate}
                \item Linux, Shell e organização de projetos
                \item Git, ambientes de desenvolvimento e organização de projetos de dados
                \item NumPy e computação numérica
                \item Pandas: manipulação e limpeza de dados
                \item Visualização de dados e análise exploratória
                \item Formatos de dados + Bancos de dados
                \item Feature Engineering e pipelines de dados
                \item Avaliação 1
            \end{enumerate}
        \end{column}
        \begin{column}{0.5\textwidth}
            \textbf{Tópicos Avançados (Semanas 9-15):}
            \begin{enumerate}
                \setcounter{enumi}{8}
                \item Web Scraping e APIs
                \item Clusterização (K-Means, DBSCAN)
                \item Busca textual e ranking
                \item Processamento de linguagem natural (NLP)
                \item Séries temporais
                \item ETL e dashboards interativos (Streamlit)
                \item Avaliação 2
            \end{enumerate}
        \end{column}
    \end{columns}
\end{frame}

\section{Por que aprender Linux?}

\begin{frame}[c]{Por que cientistas de dados usam Linux?}

\begin{itemize}

\item A maioria dos servidores roda \highlight{Linux}

\item Ferramentas de ciência de dados nasceram nesse ecossistema

\item Pipelines de dados são executados em \highlight{linha de comando}

\item Automação e reprodutibilidade

\end{itemize}

\begin{display}[Exemplos]
\begin{itemize}

\item servidores cloud

\item clusters HPC

\item pipelines ETL

\item containers Docker

\end{itemize}
\end{display}

\end{frame}

\section{História}

\begin{frame}[c]{Breve história do Unix}

\begin{itemize}

\item Década de 1960: projeto \highlight{Multics}

\item 1969: Ken Thompson e Dennis Ritchie criam o \highlight{Unix}

\item Filosofia: pequenos programas que fazem \highlight{uma coisa bem}

\item Unix se espalha pelas universidades

\end{itemize}

\end{frame}

\begin{frame}[c,plain]{Breve história do Unix}
    \centering
    \includegraphics[height=0.9\textheight]{images/Unix_history-simple.png}
\end{frame}

\begin{frame}[c]{Nascimento do Linux}

\begin{itemize}

\item 1991: \highlight{Linus Torvalds} cria o kernel Linux

\item Inspirado pelo design do Unix

\item Licença open source

\item Comunidade global de desenvolvimento

\end{itemize}

\begin{display}[Hoje]
Linux roda:

\begin{itemize}
\item servidores
\item Android
\item supercomputadores
\item infraestrutura de internet
\end{itemize}

\end{display}

\end{frame}

\section{Terminal}

\begin{frame}[c]{Interface de Linha de Comando}

\begin{columns}

\begin{column}{0.5\textwidth}


\textbf{CLI}

\begin{itemize}
\item comandos
\item scripts
\item automação
\end{itemize}\pause{}
\begin{figure}
    \centering
    \includegraphics[width=0.4\linewidth]{images/Python-logo-notext.svg.png}
\end{figure}
\pause{}

\end{column}

\begin{column}{0.5\textwidth}
\textbf{GUI}

\begin{itemize}
\item mouse
\item janelas
\item ícones
\end{itemize}\pause{}
\begin{figure}
    \centering
    \includegraphics[width=0.4\linewidth]{images/clash.jpg}
\end{figure}



\end{column}

\end{columns}

\end{frame}

\begin{frame}[c]{Comandos básicos}

\begin{columns}

\begin{column}{0.5\textwidth}
\begin{itemize}

\item ls

\item cd

\item pwd

\item mkdir

\item cp

\item mv

\item rm

\end{itemize}
\end{column}


\begin{column}{0.5\textwidth}
\begin{figure}
    \centering
    \includegraphics[width=\linewidth]{images/Televideo925Terminal.jpg}
\end{figure}
    
\end{column}

\end{columns}
\end{frame}

\section{Organização de projetos}

\begin{frame}[c, fragile]{Estrutura de um projeto de dados}

\begin{verbatim}
project/

data/
notebooks/
src/
scripts/
results/
README.md
\end{verbatim}

\begin{itemize}

\item Separar dados, código e resultados

\item Facilitar reprodução

\end{itemize}

\end{frame}

\section{Git}

\begin{frame}[c, fragile]{Versionamento}

\begin{display}[Por que usar Git?]

\begin{itemize}

\item Histórico de mudanças

\item Colaboração

\item Reprodutibilidade

\end{itemize}

\end{display}

\begin{verbatim}
git clone
git add
git commit
git push
\end{verbatim}

\end{frame}

\begin{frame}[c]{Resumo}

\begin{display}[Hoje vimos]

\begin{itemize}

\item História do Unix e Linux

\item Interface de linha de comando

\item Comandos básicos

\item Organização de projetos

\item Git

\end{itemize}

\end{display}

\end{frame}


\begin{frame}[c, noframenumbering, plain]
    \frametitle{~}
        \vfill
        \begin{center}
            {\Huge Obrigado!}\vspace{1.5em}\\
            {\Large \highlight{Alguma dúvida?}}\\
        \end{center}
        \vfill
        \begin{center}
            {\small Vamos abrir o terminal!}
        \end{center}
\end{frame}

\end{document}